\documentclass[12pt]{article}
%Gummi|065|=)
\usepackage[utf8]{inputenc}
\usepackage{graphicx}
\usepackage{wrapfig}

\title{\textbf{Načrt spletne aplikacije}}
\author{Maj Smerkol  - 63120278\\
Študijsko leto 2014/15 \\
Spletno programiranje}
\begin{document}

\maketitle

\section{Spletna aplikacija}
ToDo List \includegraphics[scale=0.1]{todo_logo.png} je spletna aplikacija, namenjena vzpodbujanju uporabnika k izpolnjevanju obveznosti. To doseže tako, da uporabniku pripravi dnevni plan in ga opozarja na pomemebna opravila.

Uporabnik doda opravila. Vsakemu opravilu nastavi oznako (ime), rok za dokončanje, pomembnost in opis, v katerem poda podrobnosti o opravilu. Namen opisa je, da uporabnik lažje prikliče v spomin opravilo.

\section{Ciljna publika}
Ciljna publika aplikacije ToDo List so srednje zahtevni uporabniki, ki želijo predvsem dostop do opravil prek interneta in opozarjanje. V tej skupini sem prepoznal predvsem študente in dijake, saj poslovni uporabniki potrebujejo več.

\section{Organizacija spletišča}
Aplikacija je organizirana naključno, vendar zelo preprosto, da se uporabnik ne izgubi na strani.

\begin{wrapfigure}{l}{0.6\textwidth}\centering
	\includegraphics[scale=1]{sitemap01.png}
\end{wrapfigure}

S strani \texttt{login} se uporabnik vpiše na stran ali pa se registrira.

Na strani \texttt{home} vidi plan za današnji dan in vse ostale naloge, ki jih mora opraviti. Tukaj lahko doda nova opravila, spreminja obstoječa in jih izbriše.

Na strani \texttt{categories} vidi vsa opravila, urejena v kategorije. Tukaj lahko opravila spreminja in jih briše.

Na strani \texttt{stats} lahko uporabnik vidi statistične podatke o svoji uspešnosti. Aplikacija izriše graf števila opravljenih nalog v odvisnosti od časa.

Izpisani so tudi nekateri drugi podatki. Ker je ciljna publika mlada, se izpišejo tudi uporabnikovi dosežki. To naj bi uporabnika motiviralo k delu.

Iz menijske vrstice lahko z vseh strani, razen \texttt{login}, dostopa do vseh ostalih strani.

\begin{wrapfigure}{l}{0.5\textwidth}\centering
	\vspace{-30pt}
	\includegraphics[width=0.5\textwidth]{wireframe_main.png}
	\vspace{-100pt}
\end{wrapfigure}

\section{Struktura spletišča}

Pri oblikovanju strukture sem se zgledoval po strani Google Keep, ki je po funkcijonalnostih prav tako podobna. 

Držal sem se klasične strukture z glavo, telesom in nogo. V glavi je logo in menijska vrstica, v telesu je vsa vsebina, v nogi pa kontaktne informacije.

\newpage
\section{Oblikovanje}
Vse strani so oblikovane po enakem principu. Glava in noga sta povsod enaki, saj to pripomore k enotnosti izgleda strani. Na vseh straneh spletišča se pojavljajo podobni elementi: telo je razdeljeno z naslovi, ki razlagajo funkcionalnost dela strani, pod njimi pa so elementi. Da se bolje loči dele strani, so naslovi napisani na kontrastnem ozadju. Elementi vsebine so povsod v poljih temnnejše barve, da se lahko posamezne elemente loči, kljub temu pa zaradi vmeščenosti med naslove delujejo kot celota.

Ciljna publika je mlada, vendar pa je ta aplikacija namenjena produktivnosti in ne zabavi, zato sem se odločil za ne preveč svetlo barvno shemo. Modra barva, na kateri celotna shema temelji, uporabnika pomirja. Celotna shema je precej hladna, s čimer nakazuje na resen namen aplikacije. 
\newline
\\
\begin{center}
\includegraphics[scale=1.6]{barvna_shema.png}
\end{center}

\end{document}
